% ====================================================================
% ФАЙЛ: tasks/03_electrodynamics/joule_lenz_law_and_power_part1.tex
% ТЕМА: РАБОТА И МОЩНОСТЬ ТОКА. ЗАКОН ДЖОУЛЯ — ЛЕНЦА
% ПАЧКА 1: ВАРИАНТЫ 1–10
% ====================================================================

% === ЗАДАЧА 1 ===
% Тег: #ЭЛЕКТРОДИНАМИКА #ТОК #МОЩНОСТЬ #ДЖОУЛЬ_ЛЕНЦ #КИМ_11 #ВАРИАНТ_06
\begin{taskbasic}[code={\label{task:elem_pow_v6_n11}}]
\textbf{Задача 1.} По проводнику сопротивлением $10\text{ Ом}$ течёт постоянный электрический ток. Величина заряда, прошедшего через поперечное сечение проводника, возрастает с течением времени согласно графику. Определите мощность, выделяющуюся в проводнике. Ответ выразите в ваттах (Вт).

\nopagebreak\smallskip
\begin{center}
\begin{tikzpicture}[scale=1.0, every node/.style={font=\footnotesize}]
    % Сетка gray!30
    \draw[xstep=1, ystep=1, gray!30, thin] (0,0) grid (4,4);
    
    % Оси координат
    \draw[-{Stealth[scale=0.8]}, black, thick] (0,0) -- (4.5,0) node[right] {$t$, с};
    \draw[-{Stealth[scale=0.8]}, black, thick] (0,0) -- (0,4.5) node[above] {$q$, Кл};
    
    % График q(t)
    \draw[ultra thick, brandPrimary] (0,0) -- (4,4);
    
    % Точки на узлах
    \fill[black] (1,1) circle (1.5pt);
    \fill[black] (2,2) circle (1.5pt);
    \fill[black] (3,3) circle (1.5pt);
    \fill[black] (4,4) circle (1.5pt);
    
    % Пунктиры
    \draw[dashed, gray] (4,0) -- (4,4);
    \draw[dashed, gray] (0,4) -- (4,4);
    
    % Оцифровка
    \node[left] at (0,1) {0,2};
    \node[left] at (0,2) {0,4};
    \node[left] at (0,3) {0,6};
    \node[left] at (0,4) {0,8};
    
    \node[below] at (1,0) {1};
    \node[below] at (2,0) {2};
    \node[below] at (3,0) {3};
    \node[below] at (4,0) {4};
    \node[below left] at (0,0) {0};
\end{tikzpicture}
\end{center}

\kim{11}
\end{taskbasic}
\answer{0,4}

% === ЗАДАЧА 2 ===
% Тег: #ЭЛЕКТРОДИНАМИКА #ТОК #МОЩНОСТЬ #ДЖОУЛЬ_ЛЕНЦ #КИМ_11 #ВАРИАНТ_09
\begin{taskbasic}[code={\label{task:elem_pow_v9_n11}}]
\textbf{Задача 2.} При протекании по проводнику постоянного электрического тока силой $0{,}5\text{ А}$ в нём за $1\text{ мин}$ выделяется количество теплоты, равное $150\text{ Дж}$. Чему равно сопротивление проводника? Ответ выразите в омах (Ом).

\kim{11}
\end{taskbasic}
\answer{10}

% === ЗАДАЧА 3 ===
% Тег: #ЭЛЕКТРОДИНАМИКА #ТОК #МОЩНОСТЬ #ДЖОУЛЬ_ЛЕНЦ #КИМ_11 #ВАРИАНТ_10
\begin{taskbasic}[code={\label{task:elem_pow_v10_n11}}]
\textbf{Задача 3.} По проводнику сопротивлением $20\text{ Ом}$ течёт постоянный электрический ток силой $0{,}5\text{ А}$. Какое количество теплоты выделится в проводнике за $1\text{ мин}$? Ответ выразите в джоулях (Дж).

\kim{11}
\end{taskbasic}
\answer{300}

% ====================================================================
% ФАЙЛ: tasks/03_electrodynamics/joule_lenz_law_and_power_part3.tex
% ТЕМА: РАБОТА И МОЩНОСТЬ ТОКА. ЗАКОН ДЖОУЛЯ — ЛЕНЦА
% ПАЧКА 3: ВАРИАНТЫ 21–30
% ====================================================================

% === ЗАДАЧА 1 ===
% Тег: #ЭЛЕКТРОДИНАМИКА #ТОК #МОЩНОСТЬ #ДЖОУЛЬ_ЛЕНЦ #КИМ_11 #ВАРИАНТ_23
\begin{taskbasic}[code={\label{task:elem_pow_v23_n11}}]
\textbf{Задача 1.} Во сколько раз уменьшится мощность электрического тока, выделяющаяся в резисторе постоянного сопротивления, если напряжение на нём уменьшить в $1{,}5$ раза?

\kim{11}
\end{taskbasic}
\answer{2,25}

% === ЗАДАЧА 2 ===
% Тег: #ЭЛЕКТРОДИНАМИКА #ТОК #МОЩНОСТЬ #ДЖОУЛЬ_ЛЕНЦ #КИМ_11 #ВАРИАНТ_24
\begin{taskbasic}[code={\label{task:elem_pow_v24_n11}}]
\textbf{Задача 2.} Во сколько раз увеличится мощность электрического тока, выделяющаяся в резисторе постоянного сопротивления, если силу тока через резистор увеличить в $2$ раза?

\kim{11}
\end{taskbasic}
\answer{4}

% === ЗАДАЧА 3 ===
% Тег: #ЭЛЕКТРОДИНАМИКА #ТОК #МОЩНОСТЬ #ДЖОУЛЬ_ЛЕНЦ #КИМ_23 #ВАРИАНТ_27
\begin{taskbasic}[code={\label{task:elem_pow_v27_n23}}]
\textbf{Задача 3.} По проводнику сопротивлением $R = 10\text{ Ом}$ течёт постоянный электрический ток силой $I = 1\text{ А}$. Какое количество теплоты выделится в проводнике за интервал времени $\Delta t = 6\text{ с}$? Ответ выразите в джоулях (Дж).

\kim{23}
\end{taskbasic}
\answer{60}

% === ЗАДАЧА 4 ===
% Тег: #ЭЛЕКТРОДИНАМИКА #ТОК #МОЩНОСТЬ #ДЖОУЛЬ_ЛЕНЦ #КИМ_25 #ВАРИАНТ_28
\begin{taskbasic}[code={\label{task:elem_pow_v28_n25}}]
\textbf{Задача 4.} Электрическая цепь состоит из источника постоянного напряжения с ЭДС $\mathcal{E} = 12\text{ В}$ и внутреннего сопротивления $r = 1\text{ Ом}$, замкнутого на два параллельно соединённых одинаковых резистора сопротивлением $R = 10\text{ Ом}$ каждый. Какое количество теплоты выделится во внешней цепи за время $t = 10\text{ с}$? Ответ выразите в джоулях (Дж).

\kim{25}
\end{taskbasic}
\answer{180}